\clearpage
\chapter*{ABSTRAK}
\phantomsection
\addcontentsline{toc}{chapter}{ABSTRAK}
\setstretch{1.5}

Perkembangan teknologi \textit{Internet of Things} (IoT) mendorong pengembangan
sistem monitoring berbasis sensor yang terhubung dengan platform \textit{cloud}
secara kontinu. Tugas besar ini mengembangkan sistem monitoring suhu berbasis
IoT cerdas menggunakan ESP32-S3 dan sensor DS18B20 yang disimulasikan pada
platform Wokwi. Data suhu dikirimkan ke ThingSpeak melalui koneksi Wi-Fi
virtual setiap 15 detik. Sistem dilengkapi dengan kecerdasan buatan berbasis
\textit{Random Forest Regressor} yang dijalankan menggunakan Python. Model AI
membaca data historis dari ThingSpeak melalui \textit{public channel},
membentuk tujuh fitur \textit{time series}, dan memprediksi suhu 1 menit ke
depan (4 langkah ke depan). Hasil
prediksi diklasifikasikan menjadi tiga kategori: dingin ($<$25\textdegree C),
normal (25\textdegree C -- 30\textdegree C), dan panas ($>$30\textdegree C).
Keluaran AI dikirim kembali ke ThingSpeak dan ditampilkan pada dashboard
bersama suhu aktual. Evaluasi model menggunakan metrik MAE, RMSE, dan R$^2$
\textit{Score}. Hasil pengujian menunjukkan sistem berhasil melakukan akuisisi
data sensor, komunikasi IoT \textit{end-to-end}, prediksi suhu berbasis AI,
dan visualisasi dashboard secara otomatis.

\vspace{0.7cm}
\noindent\textbf{Kata Kunci:} IoT cerdas, ESP32-S3, DS18B20, ThingSpeak,
\textit{Random Forest Regressor}, prediksi suhu, Wokwi.
